\section{Related Work}\label{sec:related work}

Compression is widely used in database and data-management systems to reduce storage footprint, improve I/O efficiency, and lower data-movement costs. For floating-point scientific simulation data, error-bounded lossy compression provides an effective means of reducing storage and transfer costs while maintaining user-specified control over reconstruction error. However, its effectiveness depends strongly on the organization of the underlying data. We therefore review existing compressors for structured, unstructured, and HOHS datasets.

\subsection{Structured Tensor Compression}

A large body of research has focused on error-bounded lossy compression for scientific data defined on structured grids~\cite{lindstrom2014fixed,di2016fast,tao2017significantly,liang2018error,zhao2021optimizing,liang2022sz3,gong2023mgard,huang2023cuszp}. 
These compressors typically rely on the data being stored as multi-dimensional tensors with regular spacing, enabling the efficient use of fixed neighborhood relationships and highly specialized prediction or transform schemes.
The SZ family of compressors, including SZ2~\cite{tao2017significantly}, SZ3~\cite{liang2022sz3,zhao2021optimizing}, and QoZ2.0~\cite{liu2024high}, each follows a common prediction-based compression pipeline that includes value prediction, error quantization, and lossless encoding. 
SZ2 employs local block-based regression models to predict data values, while SZ3 adopts a dimension-by-dimension prediction scheme that uses previously reconstructed neighboring values along each dimension as references. 
ZFP~\cite{lindstrom2014fixed} represents a transform-based compressor for structured tensor data, partitioning data into fixed-size blocks and applying orthogonal transforms to decorrelate values within each block. 
SPERR~\cite{li2023lossy} provides an alternative transform-based approach, applying a wavelet transform to the input data, encoding the transform coefficients using SPECK~\cite{Pearlman2004speck}, and compressing the remaining residual information using Zstandard~\cite{ycollet2021zstd}.

\subsection{Unstructured Data Compression}

In addition to studies on structured data compression, recent work has explored error-bounded lossy compression for unstructured scientific data~\cite{ajtai2002compactly,patil2018survey,salloum2018optimal,ren2024prediction,gong2024general,wu2025enabling,li2024msz}.
One line of work~\cite{wu2025enabling,ren2024prediction} seeks to improve compressibility by reorganizing data partitioning and traversal order, thereby enhancing access locality and enabling more efficient prediction. 
These approaches may also modify in-memory data layouts to reduce the cost of neighborhood access. 
Another approach~\cite{gong2024general} maps unstructured data onto regular blocks in physical space and applies regression-based predictors within each grid block.
It is important to note, however, that HOHS datasets exhibit distinct structures. While the global mesh connectivity is unstructured, data within each spectral element is organized in a locally quasi-regular manner, allowing neighborhood access within an element at a cost comparable to that of structured tensor data, despite non-uniform physical spacing and non-collinearity of points. Existing unstructured-data compressors do not explicitly exploit this element-local regularity, leaving substantial room for improvement in both compression ratio and throughput.

\subsection{Existing Compressors for HOHS Datasets}\label{sec:related work:prior work}

Existing compression methods specifically designed for HOHS datasets remain limited~\cite{marin2016spectral,liu2024parallel,zniyed2024enhanced}. 
We identify a representative prior work~\cite{marin2016spectral} that targets data produced by Nek5000~\cite{nek5000-web-page}, where each element is defined on either Gauss-Lobatto-Legendre (GLL) or Chebyshev grids. This method applies element-wise transforms, using discrete Legendre transforms (DLT) for GLL-based elements and discrete Chebyshev transforms for Chebyshev-based elements, to convert data into a coefficient space. Transformed coefficients close to zero are represented by a single bit 0, while other coefficients are prefixed with a single bit 1 and stored using a reduced-precision floating-point representation.

This approach achieves only limited compression ratios and does not provide strict error control. In practice, the reconstruction error may exceed the specified error bound unless the specified error bound is large enough. Despite these limitations, this work highlights an important opportunity: applying element-wise transforms can produce coefficient tensors with many near-zero values, which are inherently more compressible. 